\documentclass{article}

\usepackage[utf8]{inputenc}
\usepackage[T1]{fontenc}
\usepackage{helvet}
\usepackage{geometry}
\usepackage{xcolor}
\usepackage{eso-pic}
\usepackage{fancyhdr}
\usepackage{parskip}
\usepackage{microtype}
\usepackage[english, ukrainian]{babel}
\usepackage{url}
\usepackage{amsmath}
\usepackage{amssymb}
\usepackage{amsthm}
\usepackage{longtable}
\usepackage{booktabs}
\usepackage{hyperref}
\usepackage{titlesec}
\usepackage{tocloft}

\renewcommand{\cftsecfont}{\small\bfseries}
\renewcommand{\cftsubsecfont}{\footnotesize}
\renewcommand{\cftsubsubsecfont}{\scriptsize}
\renewcommand{\cftsecpagefont}{\small\bfseries}
\renewcommand{\cftsubsecpagefont}{\footnotesize}
\renewcommand{\cftsubsubsecpagefont}{\scriptsize}

\definecolor{nato}{HTML}{293757}
\definecolor{footergray}{gray}{0.65}

\geometry{a4paper,left=3cm,right=3cm,top=3cm,bottom=3cm}
\renewcommand{\familydefault}{\sfdefault}
\setcounter{secnumdepth}{3}

\titleformat{\section}
  {\color{nato}\Huge\bfseries}{\thesection.}{1em}{}
\titlespacing*{\section}{0pt}{30pt}{12pt}

\titleformat{\subsection}
  {\color{nato}\LARGE\bfseries}{\thesubsection.}{1em}{}
\titlespacing*{\subsection}{0pt}{20pt}{8pt}

\titleformat{\subsubsection}
  {\color{nato}\Large\bfseries}{\thesubsubsection.}{1em}{}
\titlespacing*{\subsubsection}{0pt}{10pt}{6pt}

\fancyhf{}
\fancyhead[R]{\small\color{footergray} 2 червня 2026}
\fancyfoot[L]{\small\color{footergray} Комплексна система захисту інформації. Цільовий профіль безпеки вищих спеціалізованих судів.}
\fancyfoot[R]{\small\color{footergray}\thepage}

\newcommand\HUGE[1]{\fontsize{40}{40}\selectfont #1}

\renewcommand{\headrulewidth}{0pt}
\renewcommand{\footrulewidth}{0.4pt}
\renewcommand{\footrule}{\color{footergray}\hrule width \textwidth height 0.4pt}

\begin{document}

\pagestyle{fancy}

\begin{titlepage}
  \AddToShipoutPictureBG*{\AtPageLowerLeft{\color{nato}\rule{\paperwidth}{\paperheight}}}
  \color{white}\thispagestyle{empty}\vspace*{4cm}\centering
  {\Huge  \textbf{Комплексна система захисту інформації} \par} \vspace{2cm}
  {\HUGE  \textbf{Цільовий профіль безпеки вищих спеціалізованих судів} \par} \vspace{2.5cm}
  \vfill
  {\large 2026 \copyright\ Інформаційні Судові Системи \par}
\end{titlepage}

\newpage

\begin{center}
  {\Huge\color{nato}\bfseries Зміст\par}
\end{center}
\vspace{40pt}
\tableofcontents

\begin{abstract}
Цільовий профіль безпеки (ЦПБ) — індивідуальний для конкретної системи підприємства. Саме на його основі створюється/модернізується КЗЗІ. Включає лише додаткові вимоги (відмінності) відносно Галузевого профілю.
\end{abstract}

\section{AU}

\subsection{ПОДІЇ АУДИТУ (AU-2)}
a. Визначити типи подій, які система може реєструвати для підтримки функції аудиту: [Призначення: типи подій, визначені організацією, які система здатна реєструвати]; b. Координувати функції аудиту безпеки з іншими організаційними підрозділами, які вимагають інформації, пов'язаної з аудитом, для посилення взаємної підтримки та допомоги у виборі типів подій, що перевіряються; c. Визначити, які типи подій підлягають аудиту: [Призначення: визначені організацією події, що підлягають аудиту (підмножина подій, що підлягають аудиту, визначених в AU-2 a.), а також частота (або ситуація, що вимагає) проведення аудиту для кожної ідентифікованої події] d. Обґрунтувати, чому типи подій, що перевіряються, вважаються достатніми для підтримки розслідувань інцидентів (постфактум), пов'язаних з безпекою та приватністю; e. Перегляньте й оновіть типи подій, вибрані для журналювання [Призначення: частота, визначена організацією].

\vspace{10pt}
{\footnotesize\bfseries
No: 1\\
Name: au\_2\_a\\
Type: string\\
Default: nil
}

{\footnotesize Типи подій, які система здатна реєструвати, визначено для підтримки функції аудиту}


{\footnotesize\bfseries
No: 2\\
Name: au\_2\_b\\
Type: string\\
Default: nil
}

{\footnotesize Функція аудиту безпеки координується з іншими підрозділами організації, які вимагають інформації, пов'язаної з аудитом, длдля посилення взаємної підтримки та допомоги у виборі типів подій, що перевіряються}


{\footnotesize\bfseries
No: 3\\
Name: au\_2\_c\_01\\
Type: string\\
Default: nil
}

{\footnotesize Типи подій (підмножина 02\_ODP[01]) визначаються для реєстрації у системі; AU-}


{\footnotesize\bfseries
No: 4\\
Name: au\_2\_c\_02\\
Type: string\\
Default: \textquotedbl{}щорічно\textquotedbl{}
}

{\footnotesize Зазначені типи подій реєструються 02\_ODP[03] частота або ситуація>; <AU-}


{\footnotesize\bfseries
No: 5\\
Name: au\_2\_d\\
Type: string\\
Default: nil
}

{\footnotesize Надається обґрунтування, чому типи подій, що перевіряються, вважаються достатніми для підтримки розслідувань інцидентів (постфактум), пов'язаних з безпекою та конфіденційністю}


{\footnotesize\bfseries
No: 6\\
Name: au\_2\_e\\
Type: string\\
Default: \textquotedbl{}щорічно\textquotedbl{}
}

{\footnotesize Переглядаються та оновлюються типи подій, вибрані для реєстрації, частота. у системі}


{\footnotesize\bfseries
No: 7\\
Name: au\_2\_odp\_01\\
Type: string\\
Default: nil
}

{\footnotesize Визначено типи подій, які система може реєструвати для підтримки функції аудиту}


{\footnotesize\bfseries
No: 8\\
Name: au\_2\_odp\_02\\
Type: string\\
Default: nil
}

{\footnotesize Визначено типи подій (підмножина AU-02\_ODP[01]) що підлягають аудиту у системі}


{\footnotesize\bfseries
No: 9\\
Name: au\_2\_odp\_03\\
Type: list\\
Default: [\textquotedbl{}login\textquotedbl{}, \textquotedbl{}logout\textquotedbl{}, \textquotedbl{}failed\_attempt\textquotedbl{}]
}

{\footnotesize Визначено частоту або ситуацію, що вимагає проведення аудиту для кожної ідентифікованої події}


{\footnotesize\bfseries
No: 10\\
Name: au\_2\_odp\_04\\
Type: string\\
Default: \textquotedbl{}щорічно\textquotedbl{}
}

{\footnotesize Частота перегляду та оновлення типів подій, обраних для журналювання}




\subsubsection{УЗАГАЛЬНЕННЯ ЗАПИСІВ ПРО АУДИТ З ДЕКІЛЬКОХ ДЖЕРЕЛ (AU-2(1))}


\vspace{10pt}
{\footnotesize\bfseries
No: 1\\
Name: au\_2\_1\_01\\
Type: list\\
Default: [\textquotedbl{}login\textquotedbl{}, \textquotedbl{}logout\textquotedbl{}, \textquotedbl{}failed\_attempt\textquotedbl{}]
}

{\footnotesize ПОДІЇ АУДИТУ - УЗАГАЛЬНЕННЯ ЗАПИСІВ ПРО АУДИТ З ДЕКІЛЬКОХ ДЖЕРЕЛ [Вилучено: Включено в АU-12]. AU-02(02) ПОДІЇ АУДИТУ - ВИБІР ПОДІЇ АУДИТУ ЗА КОМПОНЕНТАМИ [Вилучено: Включено в АU-12]}




\subsubsection{ВИБІР ПОДІЇ АУДИТУ ЗА КОМПОНЕНТАМИ (AU-2(2))}


\vspace{10pt}
\textit{Немає параметрів для цього контролю.}
\vspace{10pt}


\subsubsection{ПЕРЕГЛЯД ТА ОНОВЛЕННЯ (AU-2(3))}


\vspace{10pt}
{\footnotesize\bfseries
No: 1\\
Name: au\_2\_3\_01\\
Type: list\\
Default: [\textquotedbl{}login\textquotedbl{}, \textquotedbl{}logout\textquotedbl{}, \textquotedbl{}failed\_attempt\textquotedbl{}]
}

{\footnotesize ПОДІЇ АУДИТУ - ПЕРЕГЛЯД ТА ОНОВЛЕННЯ [Вилучено: Включено в АU-02]}




\subsubsection{ПРИВІЛЕЙОВАНІ ФУНКЦІЇ (AU-2(4))}


\vspace{10pt}
{\footnotesize\bfseries
No: 1\\
Name: au\_2\_4\_01\\
Type: list\\
Default: [\textquotedbl{}login\textquotedbl{}, \textquotedbl{}logout\textquotedbl{}, \textquotedbl{}failed\_attempt\textquotedbl{}]
}

{\footnotesize ПОДІЇ АУДИТУ - ПРИВІЛЕЙОВАНІ ФУНКЦІЇ [Вилучено: Включено в AC-06(09)]}




\section{SC}

\subsection{ДОСТУПНІСТЬ РЕСУРСІВ (SC-7)}
a. Контролювати та управляти зв'язком на зовнішньому периметрі системи та на ключових внутрішніх периметрах всередині системи. b. Реалізувати підмережі для загальнодоступних компонентів системи, які є [Вибір: фізично; логічно] відділені від внутрішніх мереж організації. c. Підключатися до зовнішніх мереж або систем тільки через керовані інтерфейси, що складаються з пристроїв захисту периметру, і розташованих відповідно до архітектури безпеки та приватності організації.

\vspace{10pt}
\textit{Немає параметрів для цього контролю.}
\vspace{10pt}


\subsubsection{ФІЗИЧНО ВІДДІЛЕНІ ПІДМЕРЕЖІ (SC-7(1))}


\vspace{10pt}
{\footnotesize\bfseries
No: 1\\
Name: sc\_7\_1\_01\\
Type: string\\
Default: nil
}

{\footnotesize ЗАХИСТ ПЕРИМЕТРА - ФІЗИЧНО ВІДДІЛЕНІ ПІДМЕРЕЖІ [Вилучено: включено до SC-7]]}




\subsubsection{ПУБЛІЧНИЙ ДОСТУП (SC-7(2))}


\vspace{10pt}
{\footnotesize\bfseries
No: 1\\
Name: sc\_7\_2\_01\\
Type: string\\
Default: nil
}

{\footnotesize ЗАХИСТ ПЕРИМЕТРА - ПУБЛІЧНИЙ ДОСТУП [Вилучено: включено до SC-7]]}




\subsubsection{ТОЧКИ ДОСТУПУ (SC-7(3))}


\vspace{10pt}
{\footnotesize\bfseries
No: 1\\
Name: sc\_7\_3\_01\\
Type: integer\\
Default: 3
}

{\footnotesize Обмежена кількість зовнішніх мережевих підключень до системи}




\subsubsection{ЗОВНІШНІ КОМУНІКАЦІЙНІ СЛУЖБИ (SC-7(4))}


\vspace{10pt}
{\footnotesize\bfseries
No: 1\\
Name: sc\_7\_4\_a\\
Type: string\\
Default: nil
}

{\footnotesize Реалізовано керований інтерфейс для кожної зовнішньої телекомунікаційної послуги}


{\footnotesize\bfseries
No: 2\\
Name: sc\_7\_4\_b\\
Type: list\\
Default: [\textquotedbl{}default\_deny\_rule\textquotedbl{}, \textquotedbl{}abac\_rule\_1\textquotedbl{}]
}

{\footnotesize Встановлена політика потоку трафіку для кожного керованого інтерфейсу}


{\footnotesize\bfseries
No: 3\\
Name: sc\_7\_4\_c\_01\\
Type: string\\
Default: nil
}

{\footnotesize Захищена конфіденційність інформації, що передається через кожен інтерфейс}


{\footnotesize\bfseries
No: 4\\
Name: sc\_7\_4\_c\_02\\
Type: string\\
Default: nil
}

{\footnotesize Захищена цілісність інформації, що передається через кожен інтерфейс}


{\footnotesize\bfseries
No: 5\\
Name: sc\_7\_4\_d\\
Type: list\\
Default: [\textquotedbl{}default\_deny\_rule\textquotedbl{}, \textquotedbl{}abac\_rule\_1\textquotedbl{}]
}

{\footnotesize Задокументовано кожен виняток з політики управління трафіком з обґрунтуванням місії або бізнес-потреби, а також тривалості такої потреби}


{\footnotesize\bfseries
No: 6\\
Name: sc\_7\_4\_e\_01\\
Type: list\\
Default: [\textquotedbl{}default\_deny\_rule\textquotedbl{}, \textquotedbl{}abac\_rule\_1\textquotedbl{}]
}

{\footnotesize Переглядаються винятки з політики потоку трафіку частота}


{\footnotesize\bfseries
No: 7\\
Name: sc\_7\_4\_e\_02\\
Type: list\\
Default: [\textquotedbl{}default\_deny\_rule\textquotedbl{}, \textquotedbl{}abac\_rule\_1\textquotedbl{}]
}

{\footnotesize Потрібно видалити винятки з політики потоку трафіку, які більше не підтримуються чітко визначеною місією або бізнеспотребою}


{\footnotesize\bfseries
No: 8\\
Name: sc\_7\_4\_f\\
Type: string\\
Default: nil
}

{\footnotesize Запобігається несанкціонований обмін управління із зовнішніми мережами}


{\footnotesize\bfseries
No: 9\\
Name: sc\_7\_4\_g\\
Type: string\\
Default: nil
}

{\footnotesize Публікується інформація, яка дозволяє віддаленим мережам виявляти несанкціонований трафік площини керування з внутрішніх мереж}


{\footnotesize\bfseries
No: 10\\
Name: sc\_7\_4\_h\\
Type: string\\
Default: nil
}

{\footnotesize Фільтрується несанкціонований трафік з зовнішніх мереж. трафіком плану}


{\footnotesize\bfseries
No: 11\\
Name: sc\_7\_4\_odp\\
Type: list\\
Default: [\textquotedbl{}default\_deny\_rule\textquotedbl{}, \textquotedbl{}abac\_rule\_1\textquotedbl{}]
}

{\footnotesize Визначено періодичність перегляду винятків з політики управління інформаційними потоками}




\subsubsection{ВІДМОВА ЗА ЗАМОВЧУВАННЯМ - ДОЗВІЛ ЗА ВИНЯТКОМ (SC-7(5))}


\vspace{10pt}
\textit{Немає параметрів для цього контролю.}
\vspace{10pt}


\subsubsection{ВІДПОВІДЬ НА РОЗПІЗНАНІ ПОМИЛКИ (SC-7(6))}


\vspace{10pt}
{\footnotesize\bfseries
No: 1\\
Name: sc\_7\_6\_01\\
Type: string\\
Default: nil
}

{\footnotesize ЗАХИСТ ПЕРИМЕТРА - ВІДПОВІДЬ НА РОЗПІЗНАНІ ПОМИЛКИ [Вилучено: включено до SC-7(18)]}




\subsubsection{ЗАПОБІГАННЯ ПОДІЛУ ТУНЕЛЮВАННЯ ДЛЯ ВІДДАЛЕНИХ ПРИСТРОЇВ (SC-7(7))}


\vspace{10pt}
{\footnotesize\bfseries
No: 1\\
Name: sc\_7\_7\_01\\
Type: string\\
Default: \textquotedbl{}автоматизований засіб моніторингу\textquotedbl{}
}

{\footnotesize Запобігається розділеному тунелюванню для віддалених пристроїв, що підключаються до систем організації, якщо розділене тунелюванню не захищено за допомогою засоби захисту}


{\footnotesize\bfseries
No: 2\\
Name: sc\_7\_7\_odp\\
Type: string\\
Default: nil
}

{\footnotesize Визначені гарантії безпечного прокладання розділеному тунелюванню}




\subsubsection{ЗАХИСТ ПЕРИМЕТРА -МАРШРУТИЗАЦІЯ АВТЕНТИФІКОВАНИХ ПРОКСІ-СЕРВЕРІВ (SC-7(8))}


\vspace{10pt}
{\footnotesize\bfseries
No: 1\\
Name: sc\_7\_8\_01\\
Type: string\\
Default: nil
}

{\footnotesize SC-07(08) \_ODP[01] внутрішній комунікаційний трафік> спрямовується до <SC-07(08) \_ODP[02] зовнішніх мереж> через автентифіковані проксі-сервери на керованих інтерфейсах}




\subsubsection{ЗАХИСТ (SC-7(9))}


\vspace{10pt}
{\footnotesize\bfseries
No: 1\\
Name: sc\_7\_9\_a\_01\\
Type: string\\
Default: nil
}

{\footnotesize Виявлено вихідний комунікаційний трафік, що становить загрозу для зовнішніх систем}


{\footnotesize\bfseries
No: 2\\
Name: sc\_7\_9\_a\_02\\
Type: string\\
Default: nil
}

{\footnotesize Заборонено вихідний комунікаційний трафік, що становить загрозу для зовнішніх систем}


{\footnotesize\bfseries
No: 3\\
Name: sc\_7\_9\_b\\
Type: string\\
Default: nil
}

{\footnotesize Перевіряється ідентичність внутрішніх пов'язаних з відмовою у зв'язку. користувачів,}




\subsubsection{ІЗОЛЬОВАНІ ВІД ІНШИХ ВНУТРІШНІХ КОМПОНЕНТІВ СИСТЕМИ ШЛЯХОМ ВПРОВАДЖЕННЯ ФІЗИЧНО ВІДОКРЕМЛЕНИХ ПІДМЕРЕЖ З КЕРОВАНИМИ ІНТЕРФЕЙСАМИ ДО ІНШИХ КОМПОНЕНТІВ СИСТЕМИ (SC-7(13))}


\vspace{10pt}
{\footnotesize\bfseries
No: 1\\
Name: sc\_7\_13\_01\\
Type: string\\
Default: \textquotedbl{}автоматизований засіб моніторингу\textquotedbl{}
}

{\footnotesize Ізольовані засоби, механізми та компоненти підтримки інформаційної безпеки від інших внутрішніх компонентів системи шляхом впровадження фізично відокремлених підмереж з керованими інтерфейсами до інших компонентів системи}


{\footnotesize\bfseries
No: 2\\
Name: sc\_7\_13\_odp\\
Type: string\\
Default: \textquotedbl{}автоматизований засіб моніторингу\textquotedbl{}
}

{\footnotesize Визначені інструменти, механізми та компоненти підтримки інформаційної безпеки, які мають бути ізольовані від інших внутрішніх компонентів системи}




\subsection{ВСТАНОВЛЕННЯ КЛЮЧАМИ (SC-12)}
Встановити та управляти криптографічними ключами для криптографічних засобів, які використовуються в системі відповідно до [Призначення: визначені організацією вимоги до генерації, поширення, зберігання, доступу та знищення ключів].

\vspace{10pt}
{\footnotesize\bfseries
No: 1\\
Name: sc\_12\_01\\
Type: string\\
Default: \textquotedbl{}AES-256-GCM\textquotedbl{}
}

{\footnotesize Встановлюються криптографічні ключі, коли в системі використовується криптографія відповідно до < SC-12\_ODP вимог >}


{\footnotesize\bfseries
No: 2\\
Name: sc\_12\_02\\
Type: string\\
Default: \textquotedbl{}AES-256-GCM\textquotedbl{}
}

{\footnotesize Здійснюється управління криптографічними ключами, коли в системі використовується криптографія, відповідно до вимог }


{\footnotesize\bfseries
No: 3\\
Name: sc\_12\_odp\\
Type: string\\
Default: nil
}

{\footnotesize Визначені вимоги до генерації, розповсюдження, зберігання, доступу та знищення ключів}




\subsection{ЗАХИСТ ІНФОРМАЦІЇ В СТАНІ СПОКОЮ (SC-28)}
Забезпечити [Вибір (один або кілька): конфіденційність; цілісність] [Призначення: визначена організацією інформація] в стані спокою.

\vspace{10pt}
{\footnotesize\bfseries
No: 1\\
Name: sc\_28\_odp\\
Type: string\\
Default: nil
}

{\footnotesize Вибрано одне або декілька з наступних ЗНАЧЕНЬ ПАРАМЕТРІВ: \{конфіденційність; цілісність\}}


{\footnotesize\bfseries
No: 2\\
Name: sc\_28\_odp\_02\\
Type: string\\
Default: nil
}

{\footnotesize Є інформація в стані спокою, яка потребує захисту}




\end{document}
